\begin{proposition}[A sufficient color-order condition]
\label{prop:color-order}
Let $G$ be a nonempty simple undirected graph and $0<k\leq h\leq2k$.
Suppose $C_1,\ldots,C_r$ partition $V(G)$ into nonempty independent sets
of $G^2$. If $h>k$, suppose additionally that no edge of $G$ joins
consecutive classes $C_i,C_{i+1}$. Then
\[
 f(v)=k(i-1)\quad\text{for }v\in C_i
\]
is an $L(h,k)$-labeling of span $k(r-1)$.
If $G^2$ also contains a clique of size $r$, this span is optimal.
\end{proposition}
\begin{proof}
Vertices in the same class have distance greater than two, so equal
labels violate no constraint. Distinct classes have label difference at
least $k$, satisfying every distance-two constraint. If $h=k$, this also
suffices for edges. Otherwise, an edge cannot join consecutive classes
by assumption, and its endpoints differ by at least $2k\geq h$.
Every class is nonempty, so the span is $k(r-1)$. A square clique of size
$r$ gives the matching lower bound by Proposition~\ref{prop:square-clique}.
\end{proof}

\paragraph{Bounded construction.}
The implementation constructs a maximum-cardinality matching in the
complement of $G^2$. Each matched pair forms a class; unmatched vertices
form singleton classes. Every class is therefore independent in $G^2$.
This restricts the construction to classes of size at most two and need
not find a minimum coloring of $G^2$. Maximum matching only minimizes
the number of classes within this restricted partition scheme; its
optimality is unnecessary for checking a resulting labeling.

For $h>k$, define a compatibility graph on the classes, with an edge
between two classes exactly when no original graph edge joins them.
The required ordering is a Hamiltonian path in this compatibility graph.
A seeded insertion heuristic starts at a low-degree class, repeatedly
inserts an unplaced class into a compatible position, and restarts if
it cannot continue. It makes at most 100 attempts, with algorithmic seed 0
and canonical vertex and class ordering. The algorithmic seed is separate
from each archived graph-generation seed. A failure is reported as
\texttt{UNKNOWN}; it proves neither nonexistence of an ordering nor
infeasibility of the labeling instance. When $h=k$, no ordering search
is needed. Parameters outside the sufficient condition are not handled
by this construction; the general SAT/ILP formulations retain their
original domain.

